%% ma: L12-B15-0011
%% lop: 12
%% chuong: 5
%% bai: 15
%% dang: 12.15.9
%% loai: TLN
%% muc: VD
%% dap_an: "2,5"
%% nguon: "(NBV)-ĐỀ SỐ 14-MỖI NGÀY 1 ĐỀ THI 2025.pdf (D14, MỖI NGÀY 1 ĐỀ - Đề số 14), Phần III Câu 4"
%% kien_thuc: "Bài 15 (phương trình đường thẳng), Bài 14 (phương trình mặt phẳng), Bài 7-8 (tọa độ trung điểm)"
%% trang_thai: chinh-thuc
%% ghi_chu: "Đã duyệt 15/10/2026. Đáp án gốc đúng (kiểm bằng sympy). Đề gốc không có hình."
\begin{ex}
Trong không gian $Oxyz$, cho điểm $A(1;-1;2)$, đường thẳng $d:\dfrac{x+1}{2}=\dfrac y1=\dfrac{z-2}{1}$ và mặt phẳng $(P):x+y-2z+5=0$. Xét đường thẳng $\Delta$ cắt $d$ và $(P)$ tại hai điểm $M$, $N$ sao cho $A$ là trung điểm của đoạn thẳng $MN$. Biết vectơ $\vec u=(1;a;b)$ là một vectơ chỉ phương của $\Delta$. Tính $a+b$.
\shortans{2,5}
\loigiai{$M\in d$ nên $M(-1+2t;t;2+t)$. $A$ là trung điểm của $MN$ nên $N(3-2t;-2-t;2-t)$.
$N\in(P)$: $(3-2t)+(-2-t)-2(2-t)+5=0\Leftrightarrow2-t=0\Leftrightarrow t=2$.
Suy ra $M(3;2;4)$, $N(-1;-4;0)$, $\overrightarrow{MN}=(-4;-6;-4)=-4\left(1;\dfrac32;1\right)$.
Vậy $a=\dfrac32$, $b=1$ và $a+b=2{,}5$.}
\end{ex}
